\begin{abstract}
  % Intro sentence
  During World War II a large scale forced migration of over 100,000
  Japanese American immigrants and their citizen children
  lead to relocation away from the West Coast.
  % Methods
  I analyse the effect that this forced relocation had
  on the post-war earnings of a linked sample of likely internees.
  My empirical strategy uses an individual-level panel of Japanese
  and Chinese men and women to compare earnings before and after internment.
  % Results
  I find that individuals who were likely interned experienced faster growth in
  earnings compared to their counterfactual earnings if they had not been
  interned.  
  % Mechanisms
  The possible mechaisms which might explain this seeming adaptation
  include occupational change and later migration decisions.
\end{abstract}